% =============================================================================
% 5.2 意图识别与参数补全评测
% 草稿版本：v0.2 (2026-05-07)
%   v0.1：作为 5.2.1 子小节起草，依附于 5.2 节（典型气象任务场景下的技术验证）。
%   v0.2：按导师意见 3 与意见 5 重构。原 5.2.1 提级为 5.2 \subsection，
%        删除原 5.2 节级导言段（与 5.1 节 5 份评测集汇总表内容重叠）。
% 字数：约 1300 字
% 风格规则：v0.2 风格（无 \textbf / \textit / 引例括号 / 双引号 / 破折号）
% 引用：\cite{ref10, ref18, ref29}
% 图表：
%   - 表 \reftab{tab:intent-overall}     35 条总表
%   - 表 \reftab{tab:intent-by-category} 按场景拆分
% 依赖宏包：booktabs / tabularx（主稿已加载）
% =============================================================================

\subsection{意图识别与参数补全评测}
\label{sec:intent-eval}

意图识别与参数补全位于 ReAct Agent 链条的最前端，承担把自然语言查询
转化为结构化 \texttt{WeatherIntent} 对象的责任 \cite{ref18}。本节在
\texttt{intent\_\bk{}bench.\bk{}jsonl} 共 35 条用例上评测两个模块的联合表现，
评测脚本为 \texttt{experiments/\bk{}eval/\bk{}run\_\bk{}intent\_\bk{}eval.\bk{}py}。

评测集覆盖 10 类场景，包含 8 类单轮场景共 28 条、多轮上下文继承类
\texttt{context\_\bk{}inference} 共 3 条与对抗鲁棒类 \texttt{adversarial}
共 2 条。8 类单轮场景包括当前天气、逐小时预报、逐日预报、预警、生活
指数、历史日报、历史小时报与出行建议，其中出行建议类共 5 条专门压测
出发地、经停地与到达地三类地点角色的标注能力。多轮场景每条用例自带
\texttt{context} 字段，用于把上一轮对话摘要拼接到 prompt 中以构成
最简形式的对话状态测试。对抗鲁棒类用例包括与气象无关的电影查询，
用于验证 \texttt{recognizer} 不被诱导触发误识别。

评测指标分为两阶段共 9 项加 1 项端到端综合。\texttt{recognizer} 阶段
有 5 项指标，分别对应意图准确率、地点集合召回率与精确率、地点角色
准确率与时间字段识别率。\texttt{completer} 阶段有 3 项指标，对应
\texttt{is\_\bk{}complete} 判定准确率、\texttt{follow\_\bk{}up} 触发准确率
与 \texttt{notes} 字段覆盖率。综合指标 \texttt{end\_\bk{}to\_\bk{}end\_\bk{}pass\_\bk{}rate}
要求 9 项全部通过才算端到端通过，是论文核心指标。35 条用例的总体
结果见 \reftab{tab:intent-overall}。

\begin{table}[H]
    \centering
    \caption{意图识别与参数补全评测 9 项指标总表}
    \label{tab:intent-overall}
    \song\wuhao
    \begin{tabular}{lr}
        \toprule
        指标 & 数值 \\
        \midrule
        端到端通过率 & 91.4\% \\
        意图准确率 & 97.1\% \\
        地点角色准确率 & 100.0\% \\
        地点集合召回率 & 96.6\% \\
        地点集合精确率 & 95.6\% \\
        时间字段识别率 & 91.7\% \\
        \texttt{is\_\bk{}complete} 判定准确率 & 100.0\% \\
        \texttt{follow\_\bk{}up} 触发准确率 & 100.0\% \\
        \texttt{notes} 字段覆盖率 & 50.0\% \\
        \bottomrule
    \end{tabular}
\end{table}

总表层面 91.4\% 的端到端通过率虽已达到工程可用水平，但单一聚合指标
容易掩盖不同场景类别下能力边界的差异化表现。本文按用户实际查询场景
将评测集进一步划分为 10 个类别独立统计，覆盖单轮直问与多轮上下文
继承两个使用维度，按类别拆分后的端到端通过率分布详见
\reftab{tab:intent-by-category}。这一分布是后续讨论中识别系统真实
能力边界的关键证据，同时也为生产环境中按场景维度做分级可观测性
提供了直接的指标基线。

8 类单轮场景共 28 条用例端到端全部通过，对抗鲁棒类 2 条全部正确
触发 \texttt{unknown} 加友好追问，出行建议类 5 条全部通过，
departure、arrival 与 waypoint 三类地点角色识别准确率均达 100\%。
唯一的 \texttt{life\_\bk{}index} 失败来自单点 \texttt{recognizer} 抖动，但
\texttt{completer} 模块触发了正确追问，业务层面用户体验仍然正确。
真正能拉开差距的是多轮上下文继承类 \texttt{context\_\bk{}inference}，
该类 3 条用例端到端通过率仅为 33.3\%，是全部 10 个类别中最低的一项，
也是本节后续失败模式分析的核心。

\begin{table}[!t]
    \centering
    \caption{意图识别评测按场景类别拆分}
    \label{tab:intent-by-category}
    \song\fontsize{9pt}{12pt}\selectfont
    \renewcommand{\arraystretch}{0.95}
    \begin{tabular}{lrr}
        \toprule
        类别 & 用例数 & 端到端通过率 \\
        \midrule
        \texttt{adversarial} & 2 & 100.0\% \\
        \texttt{current\_\bk{}weather} & 4 & 100.0\% \\
        \texttt{daily\_\bk{}forecast} & 5 & 100.0\% \\
        \texttt{historical\_\bk{}daily} & 2 & 100.0\% \\
        \texttt{historical\_\bk{}hourly} & 2 & 100.0\% \\
        \texttt{hourly\_\bk{}forecast} & 5 & 100.0\% \\
        \texttt{life\_\bk{}index} & 4 & 75.0\% \\
        \texttt{travel\_\bk{}advice} & 5 & 100.0\% \\
        \texttt{weather\_\bk{}warning} & 3 & 100.0\% \\
        \texttt{context\_\bk{}inference} & 3 & 33.3\% \\
        \bottomrule
    \end{tabular}
\end{table}

回到 \texttt{context\_\bk{}inference} 类的 33.3\% 通过率，与单轮场景
之间 66.7 个百分点的差距清晰暴露了系统在多轮长程依赖上的真实瓶颈。
失败用例的共同特征是，当用户
查询仅切换其中一个槽位时，例如从查北京明天天气切换为问那成都呢，
大语言模型能够正确识别地点切换，但会遗漏从上下文继承的时间字段
\texttt{date} 等于明天。当查询同时切换意图与隐式继承多个槽位时，
例如上一轮问武汉明天天气，本轮问那需要防晒吗，意图切换正确但
\texttt{locations} 与 \texttt{time} 两个槽位都被重置为空。这一现象
表明大语言模型在单轮意图理解上已接近人类水平，但结构化字段的跨轮
继承能力是另一个独立的能力维度。单纯把上下文摘要拼接到 prompt 中
不足以稳定实现地点继承加时间继承加意图切换的复合任务，必须借助显式
对话状态跟踪机制 \cite{ref29}。这一局限性在 \ref{sec:future-work} 节
作为未来工作展开。

分层架构的容错价值同样需要在失败用例分析中专门指出。在 \texttt{life\_\bk{}index}
类用例 \texttt{index\_\bk{}004\_\bk{}no\_\bk{}loc} 中，\texttt{recognizer} 把未提供
地点的查询识别为意图正确加地点字段填入虚构的当前位置三元组，构成
礼貌性补白幻觉。但 \texttt{completer} 模块识别出当前位置不是合法
城市名，触发了 \texttt{is\_\bk{}complete=\bk{}False} 与友好追问，最终用户体验
是系统反问您想查询哪个城市，避免了走错下游去查虚构城市的生活指数。
这一用例在严格指标下被记为失败，但业务层面是成功的。这印证了一个
论点：单点大语言模型错误并不必然导致系统失败，分层架构能够吸收单点
抖动 \cite{ref10}，\texttt{recognizer} 加 \texttt{completer} 的双阶段
设计本身就构成容错机制。
